\subsection{Infrared expansion of the coherence functional}
  \label{subsec:IR}

  This subsection addresses the constructive part of the necessity argument introduced in
  Section~\ref{sec:necessity-noninjective}.
  Having excluded injective admissible projections in
  Section~\ref{subsec:injective-holonomy},
  the goal is to show that under non-injective admissible projection
  satisfying the saturation bound of Section~\ref{sec:dynamical-equation-for-the-chi-field},
  the Einstein--Hilbert term arises as the leading infrared contribution of the coherence functional.
  This goal is not reached here.
  The subsection records what is established, where the chain stops, and which operator-level construction
  would be needed to continue it.

  To extract the dominant infrared behavior of $S_{\mathrm{coh}}$, one would expand
  $\mathcal{L}_{\mathrm{defect}}$ in powers of the curvature scale $\ell_c^{-2}$,
  where $\ell_c$ is the characteristic scale of curvature variations.
  In the infrared regime $\ell_c \gg \ell_\chi$ (with $\ell_\chi = \hbar_\chi / c_{\mathrm{BI}}$
  the microscopic capacity length), curvature is slowly varying and such an expansion would be
  controlled.

  \paragraph{Step 1: metric from the operator symbol.}
    We adopt the working identification of the effective metric $g_{\mu\nu}$ with the principal symbol of an
    effective continuum operator $\mathcal{L}_\Pi = \nabla^\mu(A^{\mu\nu}(x) \nabla_\nu)$:
    \begin{equation}
      g^{\mu\nu}(x) \propto A^{\mu\nu}(x),
      \label{eq:metric-symbol}
    \end{equation}
    where $A^{\mu\nu}$ encodes the local connectivity and stiffness of the relational substrate.
    This is an identification, not a derived result.
    The graph-Laplacian reconstruction paper~\cite{Beau2026a} defines the effective metric as the local
    quadratic form of the spectral distance, under factorizability and spectral-regularity conditions;
    it does not construct $\mathcal{L}_\Pi$ or identify its symbol with that metric.
    Appendix~\ref{sec:born-lagrangian_derivation} gives two further characterizations, the Hessian definition
    Eq.~\eqref{eq:A12-emergent-metric_v171} and the dense-limit Laplace--Beltrami operator.
    The consistency of these characterizations with Eq.~\eqref{eq:metric-symbol} is not shown.

  \paragraph{Step 2: status of the bridge from fiber curvature to the Riemann tensor.}
    The curvature 2-form $\Omega_{\mu\nu}$ of the fiber connection acts on sections of the
    adjoint bundle $\mathrm{ad}\,P$ over $\mathcal{M}$.
    Its square $\|\Omega_{\mu\nu}\|^2_{\mathfrak{g}_\Pi}$, integrated over $\mathcal{M}$,
    defines a Yang--Mills-type functional on $\mathcal{M}$.
    Earlier versions of this appendix asserted that, through a Kaluza--Klein reduction from
    $\mathcal{C}_\chi$ to $\mathcal{M}$, the fiber curvature enters the effective operator through its
    sub-principal symbol, which would contain a universal term $\tfrac{1}{6}R$ and so generate $R$ at
    two-derivative order.
    That argument is withdrawn, for three reasons.

    First, the principal symbol does not see the connection.
    For a connection Laplacian acting on sections of a bundle over $(\mathcal{M}, g)$, the principal symbol is
    $g^{\mu\nu}k_\mu k_\nu\,\mathrm{Id}$.
    Changing the bundle connection with the base metric held fixed does not change it.
    The reduction from $\mathcal{C}_\chi$ to $\mathcal{M}$ was not specified: no total-space metric, no reduction
    sector, and no identification of the reduced operator with $\mathcal{L}_\Pi$ were given.
    A construction in which the metric depends on the connection or on other fiber data is not excluded,
    but it is not provided here.

    Second, the term $\tfrac{1}{6}R$ is not a symbol.
    A degree-one symbol is homogeneous of degree one in the cotangent variable, whereas $R$ is independent of it;
    in normal coordinates and a parallel frame at a point, the connection Laplacian has no first-order term there,
    while $R$ need not vanish.
    The $\tfrac{1}{6}R$ term belongs to the heat-kernel expansion.
    For a Laplace-type operator $D = -(g^{\mu\nu}\nabla_\mu\nabla_\nu + E)$ on a closed Riemannian manifold,
    with $\nabla$ combining the Levi-Civita and bundle connections, the second heat-kernel
    coefficient is~\cite[Eq.~(4.27)]{Vassilevich2003}
    \begin{equation}
      a_2(f, D)
      =
      (4\pi)^{-n/2}\,\frac{1}{6}
      \int_{\mathcal{M}} d^n x\,\sqrt{g}\;
      \mathrm{tr}_V\bigl\{ f\,(6E + R) \bigr\},
      \label{eq:subprincipal-symbol}
    \end{equation}
    with the convention $R = 2$ on the unit sphere $S^2$, while the bundle curvature enters first in $a_4$, through
    $\tfrac{30}{360}\,\mathrm{tr}_V(f\,\Omega_{ij}\Omega_{ij})$~\cite[Eq.~(4.28)]{Vassilevich2003}.
    Curvature dependence can also be introduced separately through $E$.

    Third, the two curvatures are independent without an explicit map between them.
    A trivial line bundle with flat connection over a curved base has $\Omega = 0$ and non-zero base Riemann
    tensor; a non-flat line-bundle connection over a locally flat base gives the reverse situation.
    An identification of the type $\|\Omega\|^2 \sim R_{\mu\nu\rho\sigma}R^{\mu\nu\rho\sigma}$ therefore requires
    a construction relating the bundle connection to a metric connection, and none is supplied here.
    The conditional hypotheses~\ref{K1}--\ref{K3} of Proposition~\ref{prop:fiber-base-curvature} do not provide it:
    they relate the base Ricci tensor, the horizontal total-space Ricci tensor and a quadratic contraction of
    $\Omega$, but they do not relate $\Omega$ to the base Riemann tensor.
    For instance, a non-flat Ricci-flat base times a flat circle fiber, with the product metric and the flat
    connection, satisfies~\ref{K1}, \ref{K2} and the vacuum example of~\ref{K3}, while $\Omega = 0$ and the base
    Riemann tensor does not vanish.

    Replacing ``symbol'' by ``heat-kernel coefficient'' does not repair the derivation.
    Using Eq.~\eqref{eq:subprincipal-symbol} would require specifying the operator, its domain and signature, and
    the spectral functional (heat trace, log-determinant, or another) whose expansion defines the defect density,
    with its weights, regularization and normalization.
    It would then require showing that $S_{\mathrm{coh}}$ is that functional, and an additional scale argument to
    pass from a short-time expansion to infrared dominance.
    None of these steps is supplied here.
    The missing element is this operator-level construction; no result here shows that it cannot be given.

    Conditions (D1)--(D4) would further constrain any such construction.
    In particular, (D3) (infrared dominance of a low-derivative action) requires that
    the action be organized by the number of derivatives, with the lowest-derivative
    terms dominating in the infrared.

  \paragraph{Step 3: integration by parts and the Gauss--Bonnet identity.}
    In four dimensions, the Gauss--Bonnet combination
    $\mathcal{G} = R_{\mu\nu\rho\sigma} R^{\mu\nu\rho\sigma} - 4 R_{\mu\nu} R^{\mu\nu}
    + R^2$ is a total derivative and does not contribute to the equations of motion.
    Therefore, among all quadratic curvature invariants, only $R^2$ and
    $R_{\mu\nu} R^{\mu\nu}$ are independent at the level of the action, modulo
    boundary terms.
    Both are of order $\ell_c^{-4}$, subdominant under (D3).

    Consequently, if a construction of the kind described in Step~2 produced a local, diffeomorphism-invariant
    geometric defect density organized by derivative order, its leading infrared term would take the form
    \begin{equation}
      \mathcal{L}_{\mathrm{defect}}^{\mathrm{geom}}
      = \kappa_{\mathrm{coh}}\, R
      + \mathcal{O}(\ell_c^{-4}),
      \label{eq:defect-IR-expansion}
    \end{equation}
    possibly together with a cosmological-constant term of zero derivative order.
    The coefficient $\kappa_{\mathrm{coh}}$ is not determined by the present argument.
    The separate normalization question for the Born--Infeld candidate completion is discussed in
    Sections~\ref{subsec:stepB-BI} and~\ref{subsec:stepB-G}.
